% =============================================================
%  MechanicPro --- Technical Report
%  Template: Professional / Modern / Clean
% =============================================================

\documentclass[12pt, a4paper]{report}

% -- Encoding & Font ------------------------------------------
\usepackage[T1]{fontenc}
\usepackage[utf8]{inputenc}
\usepackage[scaled=0.92]{helvet}
\renewcommand{\familydefault}{\sfdefault}


% -- Page Layout ----------------------------------------------
\usepackage[
top=2.5cm, bottom=2.5cm,
left=3cm,  right=2.5cm
]{geometry}
% Automatic "Month Year" date (no extra package needed)
\newcommand{\monthyear}{%
	\ifcase\month\or January\or February\or March\or April\or May\or June\or
	July\or August\or September\or October\or November\or December\fi
	\space\the\year%
}

% -- Colors ---------------------------------------------------
\usepackage{xcolor}
\definecolor{Navy}{RGB}{15, 40, 75}
\definecolor{Blue}{RGB}{37, 99, 195}
\definecolor{LightBlue}{RGB}{219, 234, 254}
\definecolor{TextGray}{RGB}{90, 105, 120}
\definecolor{LightGray}{RGB}{248, 249, 251}
\definecolor{RuleGray}{RGB}{220, 225, 232}
\definecolor{StatusGreen}{RGB}{22, 163, 74}

% -- Graphics & Drawing ---------------------------------------
\usepackage{graphicx}
\usepackage{scalerel}
\usepackage{tikz}
\usetikzlibrary{calc}

% -- Tables ---------------------------------------------------
\usepackage{booktabs}
\usepackage{tabularx}
\usepackage{array}
\usepackage{longtable}

% -- Lists ----------------------------------------------------
\usepackage{enumitem}
\setlist[itemize]{leftmargin=1.5em, itemsep=2pt}
\setlist[enumerate]{leftmargin=1.5em, itemsep=2pt}

% -- Code Listings --------------------------------------------
\usepackage{listings}
\lstset{
	basicstyle=\small\ttfamily,
	keywordstyle=\color{Blue}\bfseries,
	commentstyle=\color{TextGray}\itshape,
	stringstyle=\color{StatusGreen},
	backgroundcolor=\color{LightGray},
	frame=single,
	framerule=0pt,
	rulecolor=\color{RuleGray},
	numbers=left,
	numberstyle=\tiny\color{TextGray},
	numbersep=8pt,
	breaklines=true,
	showstringspaces=false,
	tabsize=2,
	xleftmargin=12pt,
}

% -- Headers & Footers ----------------------------------------
\usepackage{fancyhdr}
\pagestyle{fancy}
\fancyhf{}
\fancyhead[L]{%
	\small\color{TextGray}\textit{MechanicPro \textemdash{} Technical Report}%
}
\fancyhead[R]{%
	\small\color{TextGray}v1.2.0%
}
\fancyfoot[L]{%
	\small\color{TextGray}Technical Report%
}
\fancyfoot[R]{%
	\small\color{TextGray}\thepage%
}
\renewcommand{\headrulewidth}{0.4pt}
\renewcommand{\footrulewidth}{0.4pt}
\renewcommand{\headrule}{%
	\color{RuleGray}\hrule width\headwidth height\headrulewidth%
}
\renewcommand{\footrule}{%
	\color{RuleGray}\hrule width\headwidth height\footrulewidth%
}

% Plain style for chapter-opening pages
\fancypagestyle{plain}{%
	\fancyhf{}
	\fancyfoot[R]{\small\color{TextGray}\thepage}
	\renewcommand{\headrulewidth}{0pt}
	\renewcommand{\footrulewidth}{0.4pt}
	\renewcommand{\footrule}{%
		\color{RuleGray}\hrule width\headwidth height\footrulewidth%
	}
}

% -- Section Formatting ---------------------------------------
\usepackage{titlesec}

\titleformat{\chapter}[block]
{\normalfont\huge\bfseries\color{Navy}}
{\color{Blue}\thechapter}{1em}{}
[\vspace{4pt}\color{Blue}\rule{\textwidth}{1.5pt}\vspace{-6pt}]

\titlespacing{\chapter}{0pt}{0pt}{20pt}

\titleformat{\section}
{\normalfont\Large\bfseries\color{Navy}}
{\color{Blue}\thesection}{1em}{}

\titleformat{\subsection}
{\normalfont\large\bfseries\color{TextGray}}
{\thesubsection}{1em}{}

\titleformat{\subsubsection}
{\normalfont\normalsize\bfseries\color{TextGray}}
{\thesubsubsection}{1em}{}

% -- TOC Styling ----------------------------------------------
\usepackage{tocloft}
\renewcommand{\cftchapfont}{\bfseries\color{Navy}}
\renewcommand{\cftchappagefont}{\bfseries\color{Navy}}
\renewcommand{\cftsecfont}{\color{TextGray}}
\renewcommand{\cftsecpagefont}{\color{TextGray}}
\setlength{\cftbeforechapskip}{6pt}

% -- Hyperlinks -----------------------------------------------
\usepackage{hyperref}
\hypersetup{
	colorlinks  = true,
	linkcolor   = Blue,
	urlcolor    = Blue,
	citecolor   = Blue,
	pdftitle    = {MechanicPro --- Technical Report},
	pdfauthor   = {Your Name},
	pdfsubject  = {Technical Report},
	pdfkeywords = {Tauri, React, TypeScript, SQLite, Desktop Application},
}

% -- Custom Commands ------------------------------------------

% Info box
\newcommand{\infobox}[1]{%
	\vspace{6pt}
	\noindent
	\colorbox{LightBlue}{%
		\begin{minipage}{\dimexpr\linewidth-2\fboxsep}%
			\vspace{4pt}
			\small #1
			\vspace{4pt}
		\end{minipage}%
	}
	\vspace{6pt}
}

% Term definition
\newcommand{\term}[1]{\textbf{\color{Blue}#1}}

% Abstract page
\newcommand{\abstractpage}[1]{%
	\clearpage
	\thispagestyle{plain}
	\vspace*{2cm}
	\begin{center}
		{\color{Navy}\Large\bfseries Abstract}\\[0.4cm]
		\color{Blue}\rule{0.3\textwidth}{1.5pt}
	\end{center}
	\vspace{0.8cm}
	\begin{quote}
		\setlength{\parskip}{6pt}
		\normalsize #1
	\end{quote}
	\vfill
	\noindent\small\color{TextGray}
	\textbf{Keywords:} Tauri, React, TypeScript, SQLite, Desktop Application,
	Mechanic Shop Management.
	\clearpage
}
\newcommand{\coverpage}{%
	\begin{titlepage}
		\begin{tikzpicture}[remember picture, overlay]
			
			%% Top bar
			\fill[Navy]
			(current page.north west)
			rectangle
			([yshift=-4cm] current page.north east);
			
			%% Accent stripe below top bar
			\fill[Blue]
			([yshift=-4cm] current page.north west)
			rectangle
			([yshift=-4.12cm] current page.north east);
			
			%% Bottom band
			\fill[LightGray]
			(current page.south west)
			rectangle
			([yshift=3.2cm] current page.south east);
			
			%% Bottom band top border
			\fill[RuleGray]
			([yshift=3.2cm] current page.south west)
			rectangle
			([yshift=3.24cm] current page.south east);
			
		\end{tikzpicture}
		
		%% -- Top bar content ----------------------
		\vspace*{1.2cm}
		\hspace{1.2cm}
		\begin{minipage}{0.8\textwidth}
			{\color{white}\footnotesize\bfseries\MakeUppercase{Technical Report}}
			\quad
			{\color{Blue!40!white}\footnotesize Technical Report}
		\end{minipage}
		
		\vspace{3.6cm}
		
		%% -- Main title block ---------------------
		\hspace{1.2cm}
		\begin{minipage}{0.82\textwidth}
			
			{\color{Navy}\Huge\bfseries\scalebox{2.0}{MechanicPro}}
			
			\vspace{0.3cm}
			{\color{TextGray}\large Management System for a Mechanic Shop}
			
			\vspace{0.6cm}
			\color{Blue}\rule{\textwidth}{1.5pt}
			\vspace{0.3cm}
			
			{\color{TextGray}\normalsize
				System Architecture, Database Design, and Technical Overview
			}
			
		\end{minipage}
		
		\vfill
		
		%% -- Metadata block -----------------------
		\hspace{1.2cm}
		\begin{minipage}{0.82\textwidth}
			\renewcommand{\arraystretch}{1.5}
			\begin{tabular}{>{\color{TextGray}\small\bfseries}l
					>{\small}l
					@{\hskip 2cm}
					>{\color{TextGray}\small\bfseries}l
					>{\small}l}
				Author    & Eduardo Freitas Fernandes      & Version   & 1.2.0      \\
				Date      & \monthyear     & Status    &
				\textcolor{StatusGreen}{\bfseries Final}        \\
				Document  & Technical Report & Language  & English    \\
			\end{tabular}
		\end{minipage}
		
		\vspace{2.8cm}
		
	\end{titlepage}
}

% =============================================================
%  DOCUMENT
% =============================================================

\begin{document}
	
	\coverpage
	
	\abstractpage{%
		This document is the technical report for \textbf{MechanicPro}, a desktop
		application designed to manage the day-to-day operations of a mechanic shop.
		The report describes the system architecture, the technology stack, the
		database design, and the build and release process.
		
		MechanicPro is built with \term{Tauri v2} and \term{Rust} for the desktop
		shell, \term{React 19} and \term{TypeScript} for the user interface, and
		\term{SQLite} for local data persistence. The application runs entirely
		offline, with no server or cloud dependency.
		
		The intended reader is a developer who wants to understand, maintain, or
		extend the system. Usage instructions are covered separately in the User
		Manual.
	}
	
	% -- Revision History -----------------------------------------
	\chapter*{Revision History}
	
	\noindent \begin{tabularx}{\textwidth}{c l X l}
		\toprule
		\textbf{Version} & \textbf{Date} & \textbf{Description} & \textbf{Author} \\
		\midrule
		1.0.0 & Jan 2026 & Initial document created           & Your Name \\
		1.1.0 & Mar 2026 & Architecture section expanded      & Your Name \\
		1.2.0 & May 2026 & Database and build sections added  & Your Name \\
		\bottomrule
	\end{tabularx}
	
	\vspace{1cm}
	\noindent\small\color{TextGray}
	This document describes version 1.2.0 of MechanicPro.
	For changes to the application itself, refer to the project \texttt{CHANGELOG.md}.
	
	\newpage
	
	% -- Table of Contents ----------------------------------------
	\tableofcontents
	
	\newpage
	
	\listoffigures
	
	\newpage
	
	\listoftables   % uncomment if you add tables
	
	% =============================================================
	%  CHAPTER 1 --- INTRODUCTION
	% =============================================================
	\chapter{Introduction}
	
	\section{Purpose}
	% Explain why this document exists and what it covers.
	
	\section{Scope}
	% What is included (and what is not). This report covers the
	% technical aspects of MechanicPro --- architecture, database,
	% and build process. It does not cover usage instructions.
	
	\section{Intended Audience}
	% Even if it's just you, state it. Good practice.
	% e.g., developers who want to understand or extend the system.
	
	\section{Document Conventions}
	% Explain any conventions used: code formatting, note boxes, etc.
	
	\infobox{\textbf{Note:} Blue boxes like this one highlight important
		information or warnings throughout the document.}
	
	% =============================================================
	%  CHAPTER 2 --- SYSTEM OVERVIEW
	% =============================================================
	\chapter{System Overview}
	
	\section{Application Purpose}
	% What does MechanicPro do? Who is it for?
	% Keep this short --- the user manual will go into detail.
	
	\section{Technology Stack}
	
	\begin{tabularx}{\textwidth}{l l X}
		\toprule
		\textbf{Layer}    & \textbf{Technology}         & \textbf{Role} \\
		\midrule
		UI Framework      & React 19 + TypeScript        & Frontend interface \\
		Build Tool        & Vite                         & Dev server and bundler \\
		Desktop Shell     & Tauri v2 (Rust)              & Native window and OS access \\
		Database          & SQLite (via Tauri SQL plugin) & Local data persistence \\
		Routing           & react-router-dom             & Page navigation \\
		Charts            & recharts                     & Dashboard visualisations \\
		PDF Generation    & jsPDF + html2canvas          & Service document export \\
		Icons             & lucide-react                 & UI icon set \\
		\bottomrule
	\end{tabularx}
	
	\section{System Requirements}
	% Minimum OS, RAM, disk space. Supported platforms (Windows,
	% macOS, Linux). Required runtimes (WebView2 on Windows, etc.)
	
	\section{High-Level Data Flow}
	% Brief description of how data moves: user action -> React ->
	% Tauri command -> Rust -> SQLite -> response back to UI.
	% A diagram here would be ideal (TikZ or an included image).
	
	% =============================================================
	%  CHAPTER 3 --- ARCHITECTURE
	% =============================================================
	\chapter{Architecture}
	
	\section{Overview}
	% Describe the two-process model: the WebView (React) and the
	% Rust core. Explain that they cannot share memory directly.
	
	\section{Frontend --- React and TypeScript}
	
	\subsection{Component Structure}
	% Describe the main pages/components and how they are organised.
	
	\subsection{State Management}
	% How is state handled? Local component state, context, etc.
	
	\subsection{Routing}
	% Which routes exist, how navigation is structured.
	
	\section{Backend --- Tauri and Rust}
	
	\subsection{Tauri Commands}
	% Explain what Tauri commands are: Rust functions exposed to
	% the frontend. List the main command groups.
	
	\subsection{Plugin Usage}
	% Tauri SQL plugin, file system plugin, dialog plugin, etc.
	
	\section{IPC Communication}
	% How the frontend calls Rust commands (invoke), how errors
	% are handled across the boundary, data serialisation (JSON).
	
	\section{Folder Structure}
	% Annotated tree of the repository.
	
	\begin{lstlisting}[language=bash, caption={Repository structure}]
		mechanic_shop/
		+--- src/               # React + TypeScript source
		|   +--- components/    # Reusable UI components
		|   +--- pages/         # Route-level page components
		|   `--- ...
		+--- src-tauri/         # Rust / Tauri backend
		|   +--- src/
		|   |   +--- main.rs    # Entry point
		|   |   `--- ...
		|   `--- tauri.conf.json
		+--- public/            # Static assets
		+--- package.json
		`--- vite.config.ts
	\end{lstlisting}
	
	% =============================================================
	%  CHAPTER 4 --- DATABASE
	% =============================================================
	\chapter{Database}
	
	\section{Overview}
	% SQLite, local file, path on each OS, no server required.
	
	\section{Schema}
	% One subsection per table. For each:
	%   - Purpose of the table
	%   - Column name, type, constraints, description
	%   - Primary key and foreign key relationships
	
	% Example table documentation:
	\subsection{clients}
	% Description of what this table stores.
	
	\begin{tabularx}{\textwidth}{l l l X}
		\toprule
		\textbf{Column} & \textbf{Type} & \textbf{Constraints} & \textbf{Description} \\
		\midrule
		id         & INTEGER & PK, AUTOINCREMENT & Unique client identifier \\
		name       & TEXT    & NOT NULL          & Full name of the client \\
		phone      & TEXT    &                   & Contact phone number \\
		email      & TEXT    &                   & Contact email address \\
		created\_at & TEXT   & NOT NULL          & ISO 8601 creation timestamp \\
		\bottomrule
	\end{tabularx}
	
	% Repeat for each table: vehicles, service_orders, etc.
	
	\section{Entity-Relationship Diagram}
	% Include an ER diagram here (TikZ or imported image).
	% \begin{figure}[h]
		%     \centering
		%     \includegraphics[width=0.85\textwidth]{figures/er_diagram.pdf}
		%     \caption{Entity-Relationship Diagram}
		%     \label{fig:er}
		% \end{figure}
	
	\section{Business Rules}
	% Rules enforced at the database level:
	% - Cascading deletes (if any)
	% - Which fields are required vs optional
	% - Any uniqueness constraints
	
	\section{CSV Format}
	% Document the import/export format precisely.
	% Column headers, data types, date formats, encoding (UTF-8).
	
	\begin{tabularx}{\textwidth}{l l X}
		\toprule
		\textbf{Column} & \textbf{Type} & \textbf{Notes} \\
		\midrule
		% Fill in for each exported entity
		\bottomrule
	\end{tabularx}
	
	% =============================================================
	%  CHAPTER 5 --- BUILD AND RELEASE
	% =============================================================
	\chapter{Build and Release}
	
	\section{Development Environment Setup}
	% Prerequisites: Node.js version, Rust toolchain version,
	% platform-specific Tauri dependencies.
	% Step-by-step to get a working dev environment.
	
	\begin{lstlisting}[language=bash, caption={Install dependencies and start dev server}]
		npm install
		npm run tauri dev
	\end{lstlisting}
	
	\section{Build Process}
	% What happens when you run the production build.
	% How Vite bundles the frontend, how Tauri compiles Rust,
	% how the two are packaged together.
	
	\begin{lstlisting}[language=bash, caption={Production build}]
		npm run tauri build
	\end{lstlisting}
	
	\section{Output Artefacts}
	% Where the binary ends up on each platform.
	% Windows: .msi or .exe installer
	% macOS: .dmg or .app
	% Linux: .deb or AppImage
	
	\section{Configuration}
	% Key settings in tauri.conf.json worth knowing:
	% app name, version, window size, permissions, plugins enabled.
	
	% =============================================================
	%  CHAPTER 6 --- KNOWN LIMITATIONS
	% =============================================================
	\chapter{Known Limitations}
	
	% Be honest. E.g.:
	% - Single-user only, no multi-user support
	% - No data encryption at rest
	% - No automatic backup mechanism
	% - No network/cloud sync
	% - No authentication layer
	
	% =============================================================
	%  APPENDICES (optional)
	% =============================================================
	\appendix
	
	\chapter{Glossary}
	% Define technical terms used in the document.
	% IPC, Tauri, WebView, SQLite, CRUD, etc.
	
	\chapter{References}
	% Tauri documentation, React docs, SQLite docs, etc.
	% Use \url{} for links.
	
	\begin{itemize}
		\item Tauri Documentation: \url{https://tauri.app/}
		\item React Documentation: \url{https://react.dev/}
		\item SQLite Documentation: \url{https://www.sqlite.org/docs.html}
	\end{itemize}
	
\end{document}